Blockieren die Räder, so rutschen sie über die Schienen, und nur die Gleitreibungskraft bremst den Zug. Mit der Normalkraft $F_\mathrm{N} = m\,g$ (alle Räder gebremst, horizontale Strecke) gilt
\begin{align}
 m\,a &= \mu_\mathrm{G}\,m\,g
\end{align}
Die Masse fällt heraus. Mit $\mu_\mathrm{G} = \muGO$ für Stahl auf Stahl (FoTaBe S.~189) ist
\begin{align}
 a &= \aBrF = \muG\cdot\gE = \aBr \approx \result{\aBrP{0}{2}}
\end{align}
Mit rollenden Rädern könnte die Haftreibung ($\mu_\mathrm{H} = \muHO$) bis $\aHaP{0}{2}$ bremsen. Blockierende Räder verlängern also den Bremsweg und schleifen zudem flache Stellen in die Räder.
